\def\biglen{20cm} % playing role of infinity (should be < .25\maxdimen)

\def\n{23} % number of random points
\def\maxxy{2} % random points are in [-\maxxy,\maxxy]x[-\maxxy,\maxxy]

\begin{figure}[!ht]
\centering
    \begin{tikzpicture}
    % generate random points
    \pgfmathsetseed{1908} % init random with the year Voronoi published his paper ;)
    \def\pts{}
    \xintFor* #1 in {\xintSeq {1}{\n}} \do{
        \pgfmathsetmacro{\ptx}{.9*\maxxy*rand} % random x in [-.9\maxxy,.9\maxxy]
        \pgfmathsetmacro{\pty}{.9*\maxxy*rand} % random y in [-.9\maxxy,.9\maxxy]
        \edef\pts{\pts, (\ptx,\pty)} % stock the random point
    }

    % draw the points and their cells
    \xintForpair #1#2 in \pts \do{
        \edef\pta{#1,#2}
        \begin{scope}
        \xintForpair \#3#4 in \pts \do{
            \edef\ptb{#3,#4}
            \ifx\pta\ptb\relax % check if (#1,#2) == (#3,#4) ?
            \tikzset{myclip/.style={}}
            \else
            \tikzset{myclip/.style={clip}}
            \fi
            \path[myclip] (#3,#4) to[half plane] (#1,#2);
        }
        \clip (-\maxxy,-\maxxy) rectangle (\maxxy,\maxxy); % last clip
        % \pgfmathsetmacro{\randhue}{rnd}
        % \definecolor{randcolor}{hsb}{\randhue,.5,1}
        \fill[lightgray] (#1,#2) circle (4*\biglen); % fill the cell with random color
        \node[dot, very thick] at (#1,#2) {}; % and draw the point
        \end{scope}
    }
    \pgfresetboundingbox
    \draw[opacity=0.3] (-\maxxy,-\maxxy) rectangle (\maxxy,\maxxy);
    \end{tikzpicture}
    \caption{Un esempio di decomposizione di Voronoi nel piano}
\end{figure}